\cotefile{157-5}
\scanpage{528}
\mspage{439}
% archive sheet number: 527

\deleted{\unsure{Passons maintenant que|Prouvons maintenant que|Posons maintenant que}}

Dans le cas où $f$ est fibrant, admettons (comme hyp. suppl au besoin) que $\deleted{\unsure{f_s|f}},\unsure{E|\mathcal E}$ satisfasse $L$ et que $f$ soit accessible, et prenons $\pi$ assez grd pour que $\deleted{\unsure{f|E_s}}E,f$ satisfassent $L_\pi$. Nous pouvons alors terminer la démonstration de 4.15.6. Il faut prouver

\[ E(\pi)\subset E_\pi \]

donc que si $x\in\mathrm{Ob}\,E_{s\pi}$, et \unsure{considérant|considérons} un système inductif dans $E$ grd devant $\pi$, de \uline{lim} $y$, on a

\[ \varinjlim\mathrm{Hom}_{E}(x,y_i)\xrightarrow{\sim}\mathrm{Hom}_{E}(x,y). \]

\deleted{\unsure{O|D}}Donc il faut prouver que tte flèche $x\xrightarrow{u}y$ se factorise, de façon ess. unique\margin{\unsure{unq|uniq}}, par un $y_i$. Mais la flèche $s\xrightarrow{\lambda=f(u)}t$ dans $C$ se factorise, de façon ess. unique, par des $\lambda_i:s\mathrel{\unsure{\relbar|\to}}t_i$. Donc on est ramené à trouver \deleted{\unsure{un|une}} \uline{$\lambda_i$-morphisme} $x\longrightarrow y_i$ pour $i$ assez grand qui factorise. Dans le cas fibrant, cela signifie qu'on cherche $x\longrightarrow\lambda_i^*(y_i)$ qui factorise $x\longrightarrow\lambda^*(y)$, or les $\lambda_i^*(y_i)$ forment un système inductif \deleted{dans} $E_s$, et sa limite \uline{dans $E_s$} est $\deleted{\unsure{s|t}}\lambda^*(y)$ (\uline{à vérifier}). Comme $x$ est de $E_{s\pi}$, l'hom donné $x\longrightarrow\lambda^*(y)$ se factorise de
